\chapter{Az opscore teljesítménymetrika}

\noindent A gráfanalitikai elemzés alapanyagát egy értelmes, szerepkörérzékeny teljesítménypontszám adja. A League of Legends meccsek nyers adatai (ölések, halálok, arany, sebzés) önmagukban nem hordoznak elegendő szemantikát: egy support 2. halála nem ugyanolyan értelmű, mint egy ADC 2. halála. Az adatgyűjtési és PUUID-normalizálási mechanizmusok a kibővített verzió 4. fejezetében kerülnek részletezésre; itt az opscore metrika belső modelljét és formuláit tárgyaljuk.

\section{A korábbi modell korlátai}

A modell fejlesztése fokozatos módszertani finomítások eredménye volt. A kiindulópontot egy nyolc-artifactos rendszer adta, amely szerepkör-szorzókat alkalmazott, majd a kapott nyers összegeket az adatbázis percentilis-eloszlásához viszonyítva normalizálta a $[0,10]$ intervallumra.

Két alapvető gyengeség hamar felszínre került. Az adatbázis-relatív normalizálás problémája: ha az adatbázis túlnyomórészt nagyon erős játékosokat tartalmaz, a medián játékos is a skála közepére eső értéket kap. A halálokon alapuló közelítés korlátai: a korábbi rendszer nem tudott különbséget tenni a tank, a support és a feeder szerepe között.

\section{A helyi, szerepkörérzékeny opscore modell}

A jelenlegi \textit{opscore} modell négy elv szerint épül fel: lokális, adatbázis-független értékelés; öt artifact tíz alkategóriával; logaritmikus bővítés az elit tartományban; megbízhatósági súlyozás rövid meccsekre.

\section{Az öt artifact és alkategóriáik}

A \ref{fig:opscore-artifacts} ábra az öt artifact és tíz alkategória teljes szerkezetét mutatja be, feltüntetve a szerepkör-érzékeny súlyozást.

\begin{sidewaysfigure}[p]
\centering
\includegraphics[width=0.95\textheight,height=0.90\textwidth,keepaspectratio]{opscore-artifacts.pdf}
\caption{Az opscore artifact-szerkezete: öt fő dimenzió, tíz alkategória, szerepkörérzékeny súlyok}
\label{fig:opscore-artifacts}
\end{sidewaysfigure}

A modell öt artifactból épül fel. A \textbf{Combat Impact} a harckonverziót és a sebzési nyomást méri, a \textbf{Risk Discipline} pedig a halálok árát és a túlélési minőséget. A \textbf{Resource Tempo} a gazdasági hatékonyságot és a korai ritmust követi. A maradék két dimenzió a csapatszintű hatást fedi le: a \textbf{Map and Objective Control} a térkép- és objektívkontrollt, a \textbf{Team Enablement} pedig a csapattársaknak teremtett értéket.

\section{Matematikai formalizmus}

\subsection{A szerepköri arány és az alapgörbe}

Minden alkategória értékelésének kiindulópontja a megfigyelt érték és a szerepköri elvárás hányadosa:
\begin{equation}
r = \frac{x}{E_{\text{role}}}
\end{equation}

Az alapgörbe egy kétparaméteres, S-alakú válaszfüggvény:
\begin{equation}
m = \left(\frac{1-t}{t}\right)^{1/h}
\end{equation}
\begin{equation}
B(r) = \frac{r^h}{r^h + m^h}
\end{equation}
ahol $t = 0{,}70$ és $h = 2$. A görbe pontosan azt biztosítja, hogy $B(1) = t$.

Az S-görbe két paramétere egy konkrét intuíciót kódol: aki pontosan az elvárásnak megfelelően teljesít ($r=1$), mindig $B(1) = 0{,}70$-et kap. Ez azt jelenti, hogy az elvárásszintű játékos a skála 70\%-án van, nem küszöb alatti, de nem is kiváló. A logaritmikus bővítő tag ($\lambda \cdot \ln r$) azokat emeli meg, akik messze felülteljesítik a szerepköri elvárást: laposabb görbe nem tudná megkülönböztetni a 120\%-on teljesítőt a 300\%-on teljesítőtől.

Pozitív mutatóknál:
\begin{equation}
S_+(r) = \min\!\left(c,\; B(r) + \lambda \cdot \ln\!\left(\max(1,\, r)\right)\right)
\end{equation}
ahol $\lambda = 0{,}50$ és $c = 2{,}20$. Negatív mutatóknál: $S_-(r) = S_+\!\left(\frac{1}{r}\right)$.

\subsection{Időtartam-megbízhatóság}

\begin{equation}
R_D = \operatorname{clamp}_{[0,1]}\!\left(\frac{D - 12}{8}\right)
\end{equation}
Ahol $D$ a meccs hossza percben. Az ingadozásra érzékeny percenkénti mutatóknál:
\begin{equation}
\tilde{S} = R_D \cdot S + (1 - R_D) \cdot 0{,}70
\end{equation}

\subsection{Artifact-aggregáció és bővítés}

\begin{equation}
A_i^{\text{alap}} = 2\,(w_1 \cdot S_1 + w_2 \cdot S_2)
\end{equation}
\begin{equation}
\gamma = 0{,}5 + 0{,}5 \cdot R_D
\end{equation}
\begin{equation}
A_i = \min\!\left(4{,}5,\; A_i^{\text{alap}} + \alpha_i \cdot \gamma \cdot \max\!\left(0,\; A_i^{\text{alap}} - 2\right)\right)
\end{equation}

Az \ref{tab:overflow_params} táblázat az öt artifact bővítési paramétereit (\(\alpha_i\)) tartalmazza.

\begin{table}[H]
\centering
\caption{Artifact-specifikus bővítési paraméterek ($\alpha_i$)}
\label{tab:overflow_params}
\begin{tabular}{|l|r|}
\hline
\textbf{Artifact} & $\boldsymbol{\alpha_i}$ \\
\hline
\textit{Combat Impact}             & $0{,}25$ \\
\textit{Risk Discipline}           & $0{,}25$ \\
\textit{Resource Tempo}            & $0{,}25$ \\
\textit{Map and Objective Control} & $0{,}35$ \\
\textit{Team Enablement}           & $0{,}20$ \\
\hline
\end{tabular}
\end{table}

\subsection{A végső opscore}

\begin{equation}
\text{opscore} = \min\!\left(10,\; \sum_{i=1}^{5} A_i\right)
\end{equation}

\subsection{Kivételes teljesítmények és szélső esetek}

A modell két jellegzetes szélső eseten keresztül szemléltethető.

\paragraph{Elité carry meccs.}
A fejlesztés során egy kiemelkedő TOP Quinn meccset vizsgáltak (kiugró ölésszám és magas assist), ahol a korábbi modell nyomás alá kerülhetett a kemény artifact-plafonok miatt. A jelenlegi modellben a \textit{Combat Impact} és a \textit{Map and Objective Control} artifactok logaritmikus bővítője lehetővé teszi, hogy ezek az értékek a nominális 2 pont fölé emelkedjenek; az összesített \textit{opscore} a 10-es plafonig futott. A halálok ára nem tűnik el: a \textit{Risk Discipline} artifact mérsékelten bünteti őket, de az elité carry összteljesítménye elismert marad.

\paragraph{Split-pusher és struktúracarry.}
Egy \textit{Yorick} főszerepű játékos, aki szándékosan kerüli a csatatereket és struktúrákon keresztül nyer, tipikusan alacsony \textit{Team Enablement} értékkel és nem kiemelkedő KDA-val rendelkezik. A \textit{Map and Objective Control} artifact dinamikus dominancia-újrasúlyozása pontosan erre az esetre van kalibrálva: ha az Objective Conversion drámaian meghaladja a Vision Control értékét, az artifact súlyt vált az objektívumszerzés felé, és az alacsony csapaterősítő szám nem torzítja indokolatlanul a végeredményt, amennyiben a win condition-beváltás ténylegesen megvalósult. A részletes szélső esetek és mintanagysági korlátok a kibővített verzióban találhatók \cite{premadegraph-full}.

\section{Az opscore szerepe a rendszerben}

Az \textit{opscore} és \textit{feedscore} pontszámok két helyen jelennek meg a rendszerben. Az asszortativitás- és centralitásszámítások csomóponti attribútumként használják ezeket: a kérdés az, hogy az összekapcsolt játékosok hasonló teljesítményprofilokat hordoznak-e. A Tribal NeuroSim szimulációban a klaszter-szintű aggregált pontszámok (\textit{A\_combat}, \textit{A\_resource}, \textit{A\_map\_objective}, \textit{A\_risk}, \textit{A\_team}) közvetlenül az öt artifactból számolódnak és a törzsek iniciális genomszekvenciáját alkotják; a meccsadatokból kiszámolt teljesítményprofil vezérli a szimulációban megjelenő viselkedési tendenciát.
